\subsection{Teoría de Números: Regla de Divisibilidad Alternada}

\graybold{Preguntas guía:}

\begin{itemize}
\item ¿El problema pide verificar o forzar divisibilidad por (base ± 1) de un número representado por dígitos?
\item ¿Se puede aprovechar que base ≡ −1 (mod base+1), lo que convierte la suma ponderada de dígitos en una suma ALTERNADA módulo (base+1)?
\item ¿Se pide además construir/ajustar un dígito específico para lograr esa divisibilidad?
\end{itemize}

\graybold{Utilidad:} Generalizar la clásica regla de divisibilidad por 11 en base 10 (suma alterna de dígitos) a cualquier base B, para verificar o forzar divisibilidad por B+1 sin operar con el número completo (que puede tener hasta 2·10\textasciicircum{}5 dígitos).

\graybold{Complejidad:} O(L), con L = cantidad de dígitos.

\graybold{Fundamento teórico:} Como B ≡ −1 (mod B+1), cada potencia B\textasciicircum{}i ≡ (−1)\textasciicircum{}i (mod B+1). Por lo tanto, el valor del número módulo B+1 es exactamente la suma alternada de sus dígitos (sumando los de posición par, restando los de posición impar). Para forzar que esa suma sea 0 módulo B+1 cambiando un solo dígito D\_i por un valor MENOR, se despeja cuánto debe valer ese dígito para anular el residuo actual, y se verifica que el nuevo valor sea válido (entre 0 y el dígito original, sin aumentarlo).

\graybold{Ejercicio aplicado}

Dado un número N de L dígitos en base B, encontrar un dígito que se pueda reducir (nunca aumentar) para que el número resultante sea múltiplo de B+1, priorizando el menor valor final posible, o determinar que no existe solución.

\graybold{I/O:} L ≤ 2·10\textasciicircum{}5 → aunque el código usa input() estándar por simplicidad (ya venía así y funciona), para esta escala también sería válido leer con sys.stdin.buffer.read().split() (patrón 2) si se buscara el máximo rendimiento.

\graybold{Código solución}

\begin{lstlisting}[language=Python]
# Lee la base B y la cantidad de digitos L
basedig = list(map(int, input().split()))
nums = list(map(int, input().split()))   # digitos D_1..D_L, de mas a menos significativo
 
base = basedig[0]
div = basedig[0] + 1                     # queremos que N sea multiplo de base+1
dig = basedig[1]                         # L = cantidad de digitos
 
# suma alternada de digitos == N mod (base+1), porque base = -1 mod (base+1)
n = 0
for i in range(len(nums)):
    if (dig - 1 - i) % 2 == 0:
        n += nums[i]
    else:
        n -= nums[i]
n %= div
 
if n == 0:
    print("0 0")                         # ya es multiplo, no hace falta cambiar nada
    exit()
 
# buscar el primer digito (de izquierda a derecha) que se pueda REDUCIR
# para anular el residuo n, segun la paridad de su posicion
res = 0
for i in range(len(nums)):
    if (dig - 1 - i) % 2 == 0:
        res = n                          # esta posicion suma: hay que restarle 'n'
    else:
        res = div - n                    # esta posicion resta: se compensa distinto
    if 1 <= res <= nums[i]:              # solo vale si es una reduccion real (no negativa)
        print(f"{i + 1} {nums[i] - res}")
        exit()
 
print("-1 -1")                           # ningun digito puede arreglarlo
\end{lstlisting}